%------------------------------------------------------------------------------
\subsection{Part II: The Enrichment Landscape}
%------------------------------------------------------------------------------

\subsubsection{Two Organizational Regimes}
\label{sec:two-regimes}

Across 9 datasets and 2,097 subjects, the $\phisym$-lattice enrichment landscape resolves into two mutually exclusive organizational regimes. The first, \textit{Alpha Mountain}, is defined by enrichment concentrated in the central octave: the attractor position ($u = 0.500$, $\approx$9.67 Hz) reaches $+41\%$, Noble$_1$ ($u = 0.618$, $\approx$10.23 Hz) reaches $+30\%$, and inv\_noble$_1$ ($u = 0.382$, $\approx$9.13 Hz) reaches $+36\%$. Both octave boundaries are depleted (lower boundary $-35\%$, upper boundary $-73\%$). This mountain profile is the empirical realization of individual alpha frequency (IAF) clustering, reflecting thalamocortical T-type Ca$^{2+}$ channel dynamics with an intrinsic $\approx$100 ms time constant that pins oscillatory peaks to the mid-octave region independent of cortical frequency-gradient forces.

The second regime, \textit{Edge Ramp}, governs theta, beta-low, and gamma, in which peaks deplete across the lower octave and accumulate monotonically toward the upper octave boundary. In beta-low, the ramp spans from noble$_4$ ($-54\%$) to inv\_noble$_6$ ($+87\%$). In theta, the ramp is the steepest of any band, rising from the lower boundary ($-52\%$) to boundary\_hi ($+123\%$)---the latter representing $\fzero$ convergence under resting conditions. Beta-high exhibits the mirror pattern: a descending ramp from its lower boundary ($+85\%$) to boundary\_hi ($-18\%$), reflecting the same $\approx$20 Hz accumulation zone viewed from the opposing octave.

Profile consistency, assessed as the fraction of positions that agree in sign direction across all 9 datasets, confirms that these regimes are robust: theta 12/13, alpha 10/13, beta-low 12/13, beta-high 4/13, and gamma 6/13 positions are consistent. The weaker consistency of beta-high under FOOOF likely reflects sensitivity of the parametric aperiodic model in this frequency range: replication with IRASA, a non-parametric aperiodic decomposition method \parencite{wen2016irasa,gerster2022irasa}, yields 13/13 cross-dataset consistency for beta-high with a steeper gradient (192 versus 113 percentage points; see Section~\ref{sec:irasa}). Gamma consistency remains weak under both methods, likely reflecting electromyographic and microsaccade artifact contamination that scalp EEG cannot fully resolve (gamma results throughout this paper should be interpreted with this caveat).

\begin{figure}[htbp]
\centering
\includegraphics[width=\textwidth]{fig4_enrichment_landscape.png}
\caption{\textbf{Enrichment landscape across 5 frequency bands and 13 $\phisym$-lattice positions.} Values are means across 9 datasets ($N = 2{,}097$ subjects total, 4.57 million peaks; Hz-corrected Voronoi enrichment, top-50\% power filter). Green indicates enrichment (more peaks than expected under uniform log-frequency distribution); red indicates depletion. Alpha shows a central mountain (attractor $+43\%$); theta, beta-low, and gamma show ascending ramps; beta-high shows a descending ramp. The alpha profile is the approximate inverse of all other bands (Pearson $r \approx -0.95$ with theta across 13 positions).}
\label{fig:enrichment_landscape}
\end{figure}

\subsubsection{Cross-Boundary Architecture}
\label{sec:cross_boundaries}

The four $\phisym$-lattice band boundaries---located at $\fzero$, $\fzero \times \phisym$, $\fzero \times \phisym^2$, and $\fzero \times \phisym^3$---exhibit qualitatively distinct enrichment signatures (Table~\ref{tab:boundaries}). Each boundary type reflects a different interaction between the intrinsic dynamics of its flanking bands.

\begin{table}[h]
\centering
\caption{Cross-boundary enrichment architecture. Enrichment values are 9-dataset means, Hz-corrected, top-50\% power filter. $\Delta$ is the signed difference (above minus below). Boundary types: Cliff (sharp drop), Void (depleted both sides), Bridge (enriched both sides), Weak (near-null both sides).}
\label{tab:boundaries}
\begin{tabular}{lrrrrrl}
\toprule
\textbf{Boundary} & \textbf{Hz} & \textbf{$\phisym$-position} & \textbf{Below} & \textbf{Above} & \textbf{$\Delta$ (pp)} & \textbf{Type} \\
\midrule
$\theta / \alpha$       & 7.60  & $\fzero$              & $+123\%$ & $-35\%$ & $-158$ & Cliff  \\
$\alpha / \beta_{\mathrm{L}}$ & 12.30 & $\fzero \times \phisym$   & $-73\%$  & $-31\%$ & $+42$  & Void   \\
$\beta_{\mathrm{L}} / \beta_{\mathrm{H}}$ & 19.90 & $\fzero \times \phisym^2$ & $+83\%$  & $+85\%$ & $+2$   & Bridge \\
$\beta_{\mathrm{H}} / \gamma$ & 32.19 & $\fzero \times \phisym^3$ & $-18\%$  & $-14\%$ & $+4$   & Weak   \\
\bottomrule
\end{tabular}
\end{table}

The $\fzero$ boundary (7.60 Hz) is the sharpest spectral transition in the dataset: theta peaks accumulate at $+123\%$ immediately below while alpha begins at $-35\%$ immediately above, a discontinuity of $-158$ pp. The alpha/$\beta_{\mathrm{L}}$ boundary at 12.30 Hz is a spectral void, with both the upper flank of alpha ($-73\%$) and the lower flank of beta-low ($-31\%$) deeply depleted---a spectral desert spanning $\approx$11.5--13.2 Hz. The $\beta_{\mathrm{H}}$/$\gamma$ boundary at 32.19 Hz is functionally inert: both flanks are near zero ($-18\%$ below, $-14\%$ above). The $\beta_{\mathrm{L}}/\beta_{\mathrm{H}}$ boundary at 19.90 Hz is the sole exception to this pattern and is described in detail in Section~\ref{sec:bridge}.

\begin{figure}[htbp]
\centering
\includegraphics[width=\textwidth]{fig5_cross_boundary.png}
\caption{\textbf{Cross-boundary architecture.} Enrichment profiles near each of the four $\phisym$-lattice band boundaries, computed as 9-dataset mean Hz-corrected Voronoi enrichment ($N = 2{,}097$ subjects). Dotted vertical lines mark the boundary frequency. (Left to right:) The $\theta/\alpha$ boundary at 7.60~Hz is a cliff ($\Delta = -161$~pp); the $\alpha/\beta_L$ boundary at 12.30~Hz is a void (both sides depleted); the $\beta_L/\beta_H$ boundary at 19.90~Hz is a bridge (both sides enriched, $\Delta = +8$~pp); the $\beta_H/\gamma$ boundary at 32.19~Hz is weak. $\Delta$ is the signed difference in enrichment (above minus below boundary).}
\label{fig:cross_boundary}
\end{figure}

\subsubsection{Alpha is the Near-Exact Inverse of All Other Bands}
\label{sec:alpha-inverse}

The most structurally distinctive feature of the enrichment landscape is the systematic sign-inversion of alpha relative to all other bands. Across all 13 $\phisym$-lattice positions, the sign of alpha enrichment is opposite to theta, beta-low, and gamma, with a position-level correlation between the alpha and theta profiles of $r \approx -0.95$.

This anti-correlation organizes the octave into three functional zones. In the lower zone ($u = 0.00$--$0.38$), alpha is depleted ($-3\%$ to $-35\%$) while beta-high is enriched and theta, beta-low, and gamma are also largely depleted. In the mid zone ($u = 0.38$--$0.62$), alpha is strongly enriched ($+30\%$ to $+41\%$) while all other bands are depleted or null. In the upper zone ($u = 0.62$--$1.00$), alpha collapses strongly ($-5\%$ to $-73\%$) while theta ($+35\%$ to $+123\%$), beta-low ($+29\%$ to $+83\%$), and gamma ($+22\%$ to $+55\%$) rise in concert.

The extreme of the anti-correlation is at boundary\_hi (12.30 Hz): theta $+123\%$ versus alpha $-73\%$, a 196 pp gap between co-occurring peaks in the same $\phisym$-octave coordinate. The biophysical interpretation is that thalamocortical T-type Ca$^{2+}$ channel dynamics provide sufficient intrinsic resonance to pin alpha peaks to the mid-octave attractor and Noble$_1$ positions, actively opposing the upper-octave gradient that governs theta, beta-low, and gamma. Where alpha occupies the mid-octave, competing oscillatory modes are suppressed; where alpha vacates the upper octave, gradient-driven accumulation is unopposed.

\subsubsection{The $\approx$20 Hz Bridge}
\label{sec:bridge}

The $\beta_{\mathrm{L}}/\beta_{\mathrm{H}}$ boundary at 19.90 Hz ($= \fzero \times \phisym^2$) is the only band boundary in the $\phisym$-lattice that is enriched from both sides simultaneously. The beta-low ascending ramp peaks at inv\_noble$_6$ ($+87\%$, 19.37 Hz) before reaching boundary\_hi ($+83\%$, 19.90 Hz), while beta-high begins its descending ramp at its lower boundary ($+85\%$, 19.90 Hz), declining monotonically to boundary\_hi ($-18\%$, 32.19 Hz). The discontinuity across the 19.90 Hz boundary is only $+2$ pp---the smallest of any boundary and the only one consistent with continuity of spectral density.

This convergence from both sides constitutes a spectral attractor at $\approx$20 Hz that is independent of paradigm, age group, and recording system: all 9 datasets (4 adult, 5 pediatric HBN releases) show enrichment at the beta-low upper positions, and the bridge pattern is confirmed in resting-state datasets without any motor component.

The 19.90 Hz frequency satisfies the Fibonacci identity $\fzero + \fzero \times \phisym = \fzero \times \phisym^2$ (i.e., $7.60 + 12.30 = 19.90$~Hz). However, this identity is a mathematical tautology---it follows from $\phisym^2 = \phisym + 1$ and holds for any $\fzero$. It does not, by itself, explain why peaks converge at this frequency. The explanation must be biophysical: four mechanisms are known to peak at $\approx$20 Hz: (1) M-current KCNQ channel resonance ($\approx$50 ms time constant); (2) corticomuscular coherence (17--21 Hz during sustained contraction); (3) post-movement beta rebound ($\approx$19.4 Hz); and (4) the antikinetic frequency (20 Hz tACS reliably slows voluntary movement). These four are not necessarily independent---M-current resonance could be the common biophysical substrate underlying corticomuscular coherence and the antikinetic effect, in which case the apparent convergence of distinct mechanisms would reduce to a single channel-level time constant expressed across motor-system contexts. Whether the convergence of these processes at a frequency that happens to equal $\fzero \times \phisym^2$ reflects a deep organising principle or a coincidence of biophysical time constants is an open question that the present data cannot resolve. The bridge is an empirical finding; its interpretation remains speculative.
